\section{Methodology}
\label{sec:method}

This section describes our experimental setup, dataset, prompt designs, evaluation metrics, statistical analysis plan, and replication package.

\subsection{Dataset}
We extracted a dataset of $N=500$ User Stories from the benchmark dataset published by Rathnayake et al.~\cite{rathnayake2026}. Each sample contains:
\begin{itemize}
    \item A User Story in Connextra format (\textit{As a... I want to... So that...}).
    \item Detailed business requirements and acceptance criteria.
    \item An expert-written Ground Truth Gherkin scenario.
\end{itemize}
Evaluation was conducted across all 500 samples using a fixed random seed of 42 to ensure reproducibility.

\subsection{Experimental Pipeline}
Our automated experimental execution pipeline consists of six sequential stages:
\begin{enumerate}
    \item \textbf{Dataset Ingestion}: Ingest $N=500$ requirements and format prompt inputs.
    \item \textbf{Prompt Formatting}: Construct Zero-Shot, Few-Shot, and Chain-of-Thought prompt instances.
    \item \textbf{GPT-4o Inference}: Execute cloud-based API inference using \texttt{GPT-4o} via the OpenAI Responses API (\texttt{temperature = 0.0}, \texttt{top\_p = 1.0}).
    \item \textbf{BDD Artifact Parsing}: Extract Gherkin scenario blocks and Python Behave step definition blocks using structural regex parsers.
    \item \textbf{Metric Computation}: Compute Sentence-BERT Cosine Similarity, Gherkin syntax validation, Python AST Syntax Parse Rates, and structural content statistics.
    \item \textbf{Statistical Analysis}: Perform One-Sample/Pairwise Wilcoxon Signed-Rank tests, Binomial exact tests, and McNemar tests to evaluate hypotheses.
\end{enumerate}

\subsection{Prompt Strategies}
We systematically evaluated three prompt engineering strategies:
\begin{itemize}
    \item \textbf{Zero-Shot Prompting}: The model was provided with the User Story and acceptance criteria and directly instructed to output the Gherkin scenario and Behave Python step definitions without any exemplars.
    \item \textbf{Few-Shot Prompting}: The prompt included two concrete in-context exemplars illustrating correct Gherkin keyword formatting (\texttt{Given-When-Then}) and Behave Python step decorators (\texttt{@given}, \texttt{@when}, \texttt{@then}).
    \item \textbf{Chain-of-Thought (CoT) Prompting}: The prompt explicitly instructed the model to generate intermediate logical reasoning steps prior to emitting Gherkin scenarios and Python step definitions.
\end{itemize}

\subsection{Evaluation Metrics and Thresholds}
We evaluated model outputs across four primary metrics:
\begin{itemize}
    \item \textbf{Sentence-BERT Cosine Similarity (Semantic Fidelity)}: Extracted Gherkin scenarios were compared against expert Ground Truth scenarios using Sentence-BERT (\texttt{all-MiniLM-L6-v2}) embedding cosine similarity. The quality baseline threshold was set to $\ge 0.80$.
    \item \textbf{Gherkin Syntax Validation}: Verified the structural presence of mandatory BDD keywords (\texttt{Feature}, \texttt{Scenario}, \texttt{Given}, \texttt{When}, \texttt{Then}).
    \item \textbf{AST Parse Rate (Syntactic Validity)}: Extracted Behave Python step definitions were evaluated using Python's Abstract Syntax Tree parser (\texttt{ast.parse()}). Code blocks compiling without \texttt{SyntaxError} received a score of 1, and 0 otherwise. The baseline threshold was set to $\ge 85\%$.
    \item \textbf{Content Statistics}: Measured structural properties including average step count, average scenario count, percentage of \texttt{And} conjunction usage, average word count, average line count, and step bounds.
\end{itemize}

\subsection{Statistical Analysis Plan}
All statistical hypothesis tests were evaluated at $\alpha = 0.05$:
\begin{itemize}
    \item \textbf{One-Sample Wilcoxon Signed-Rank Test}: Selected because Cosine Similarity scores are continuous, bounded $[0, 1]$, and non-normally distributed. Evaluates whether median Cosine Similarity for each technique significantly exceeds the $0.80$ baseline threshold.
    \item \textbf{One-Sample Binomial Exact Test}: Selected because AST parse status is binary ($1/0$). Evaluates whether the syntactic pass rate significantly exceeds the $85\%$ benchmark.
    \item \textbf{Pairwise Wilcoxon Signed-Rank Test}: Selected for paired comparison of continuous semantic similarity scores across techniques on the same 500 requirement items.
    \item \textbf{McNemar's Test}: Selected for paired comparison of binary AST parse success/failure outcomes across techniques.
    \item \textbf{Effect Size Quantification}: Calculated using Cohen's $d$ for paired continuous metrics to quantify practical difference magnitude ($|d| < 0.2$ Negligible, $0.2 \le |d| < 0.5$ Small, $0.5 \le |d| < 0.8$ Medium, $|d| \ge 0.8$ Large).
\end{itemize}

\subsection{Replication Package}
All prompts, raw API response logs, generated datasets, evaluation scripts, and statistical notebooks are available in our public GitHub repository: \url{https://github.com/Lorden13/RT-SWT-001-Gr2}.
